\documentclass[10pt,a4paper]{amsart}
\usepackage[margin=19mm]{geometry}
\usepackage{amsmath,amssymb,mathtools}
\usepackage[scr=rsfs,cal=euler]{mathalfa}
\usepackage{xfrac}
\usepackage[unicode,hidelinks,bookmarks=false]{hyperref}
\newcommand{\ie}{i.e.\ }
\newcommand{\cf}{cf.\ }
\newcommand{\resp}{respectively\ }
\newcommand{\Subsection}{\subsection}
\providecommand{\nobreakdash}{\nobreak\hbox{-}}
\setlength{\emergencystretch}{2em}
\multlinegap0pt
\usepackage[shortlabels]{enumitem}
\usepackage{amssymb}
\usepackage{mathtools}
\newtheorem{lemma}{Lemma}
\newtheorem{proposition}{Proposition}
\newcommand{\Bl}{\operatorname{Bl}}
\newcommand{\CC}{\mathbf C}
\newcommand{\OO}{\mathcal O}
\newcommand{\PP}{\mathbf P}
\newcommand{\cA}{\mathcal A}
\title{Smooth Counterexamples to the Eisenbud--Schreyer--Weyman Ulrich Existence Problem}
\author{Cristian Anghel}
\date{Source: arXiv:2609.02718v1 (CC BY 4.0)}
\begin{document}
\maketitle

\noindent\emph{Source and licence.} Smooth Counterexamples to the Eisenbud--Schreyer--Weyman Ulrich Existence Problem. Version arXiv:2609.02718v1 (CC BY 4.0), distributed by arXiv under the Creative Commons Attribution 4.0 International licence, http://creativecommons.org/licenses/by/4.0/, as recorded in the arXiv licence metadata for this exact version and shown on the abstract page https://arxiv.org/abs/2609.02718v1. Full licence notice: \texttt{licenses/arxiv-2609.02718-CC-BY-4.0.txt}.\par\smallskip

\noindent\textit{Editorial scope.} Counted unit: the article's proposition prop:local (source0.tex lines 565--599). The article's complete original proof of it is delivered with it (source0.tex lines 601--709, one \texttt{proof} environment of 109 lines). No mathematics is written, repaired or re-derived: the statement and the proof are the article's own lines, in the article's own notation, and no retained line is substituted or re-flowed.  Retained uncounted as the source-local context the proof needs: lines 454--457; lines 520--537.  Nothing was omitted from any retained span, so the package carries no omission record.  No retained line contains a citation, so the wrapper carries no bibliography and no bibliography entry.  No retained line contains a figure environment, an image-inclusion command or drawing code, so no image is delivered and the package declares no asset dependency.  Editorial formatting only, in addition: an AMS \texttt{amsart} preamble that restates the article's own declarations and macros used by the retained lines, namely the environments \texttt{lemma}, \texttt{proposition} on separate counters, together with the standard packages those lines need.  The retained environments print in this document as Lemma 1, Proposition 1 in order of appearance, whereas the article prints the same items with the numbering of its own full document.  The complete native source of the article as distributed in the arXiv e-print for this exact version is bundled unchanged as \texttt{sources/209231/source0.tex}.
\begin{lemma}\label{lem:s1-generic-field}
Let $Z$ be a Noetherian scheme satisfying $S_1$.  If $Z$ has a unique minimal
point and its local ring at that point is a field, then $Z$ is integral.
\end{lemma}

\begin{lemma}[Blowing up a completed standard cone]\label{lem:completed-cone-blowup}
Let $(R,\mathfrak m)$ be a Noetherian local $\CC$-algebra.  Suppose that its
$\mathfrak m$-adic completion is isomorphic to the completion at the vertex of a
standard graded affine cone
\[
        A=\CC[x_1,\ldots,x_r]/I,
        \qquad C=\operatorname{Proj}A,
\]
where $A$ is a normal domain and $C$ is smooth, and suppose that the isomorphism carries
$\mathfrak m\widehat R$ to the vertex ideal $(x_1,\ldots,x_r)\widehat A$.
Then $\Bl_{\mathfrak m}\operatorname{Spec}R$ is regular along its exceptional
fiber.  Scheme-theoretically that fiber is a reduced copy of $C$, and, if it is
denoted by $E$, then
\[
        \OO(-E)|_E\simeq\OO_C(1).
\]
In particular $N_{E/\Bl_{\mathfrak m}\operatorname{Spec}R}\simeq\OO_C(-1)$.
\end{lemma}

\begin{proposition}\label{prop:local}
Let $p$ be an $r$-fold point of $\cA$.
\begin{enumerate}[label=\textup{(\roman*)}]
\item The fiber of $\pi_n$ over $p$ has $n^{k-r-1}$ points.
\item If $r\le2$, $X_n$ is smooth over $p$.
\item If $r\ge3$, write $T_p=\{i:L_i\ni p\}$.  At every
      $q\in\pi_n^{-1}(p)$ there is an isomorphism of complete local
      $\CC$-algebras
      \begin{equation}\label{eq:completed-local-cone}
      \widehat\OO_{X_n,q}\simeq
      \CC[[x_i:i\in T_p]]/(F_1,\ldots,F_{r-2}),
      \end{equation}
      where each $F_a$ is a homogeneous form
      $F_a=\sum_{i\in T_p}\lambda_i^{(a)}x_i^n$, and the ambient coordinate
      $z_i$ maps to $x_i$ for every $i\in T_p$.  The right-hand side is the
      completed local ring at the vertex of the affine cone over the smooth
      complete-intersection curve
      \[
        C_{n,r}=\Phi_n^{-1}(\PP^1)\subset\PP^{r-1},
      \]
      cut out by $r-2$ hypersurfaces of degree $n$.  This curve has
      \[
        d_{n,r}:=\deg C_{n,r}=n^{r-2}
      \]
      and
      \begin{equation}\label{eq:genus}
        g_{n,r}=1+\frac12n^{r-2}\bigl((r-2)n-r\bigr).
      \end{equation}
\item Blowing up $q$ resolves the singularity in one step.  The exceptional
      curve is a reduced copy of $C_{n,r}$ with self-intersection
      \[
        C_{n,r}^2=-n^{r-2}.
      \]
\end{enumerate}
\end{proposition}

\begin{proof}
At $p$ exactly $k-r$ of the $k$ homogeneous coordinates are non-zero.  The
fiber of the power map is therefore a torsor under
$\mu_n^{k-r}/\mu_n$, giving $n^{k-r-1}$ points.

We first separate the cases with fewer than two vanishing arrangement coordinates.
If $r=0$, the power map is etale at every point of the fiber, so the completed
local ring is regular.  If $r=1$, choose local coordinates $(u,v)$ on the base
with the unique arrangement line given by $u=0$.  After extracting $n$-th roots
of the remaining units, each completed local branch has the form $t^n=u$, hence
is $\CC[[t,v]]$.  Thus $X_n$ is smooth over all points with $r\le1$.

Assume now $r\ge2$ and put $U=\{1,\ldots,k\}\setminus T_p$.
Let
\[
 W=\{(\lambda_1,\ldots,\lambda_k)\in\CC^k:
       \textstyle\sum_i\lambda_i\ell_i=0\}
\]
be the $(k-3)$-dimensional relation space.  The subspace $W_{T_p}$ of relations
supported on $T_p$ has dimension $r-2$, because the $r$ forms vanishing at $p$
span the two-dimensional space of linear forms vanishing at $p$.  Restriction
to the unit coordinates gives a linear map
\[
        \rho_U:W\longrightarrow\CC^U
\]
with kernel exactly $W_{T_p}$.  Evaluating a relation at $p$ shows that its image is
contained in the hyperplane
\[
 H_p=\Bigl\{(a_j)_{j\in U}:\sum_{j\in U}a_j\ell_j(p)=0\Bigr\}.
\]
Both $W/W_{T_p}$ and $H_p$ have dimension $k-r-1$, hence the induced map
$W/W_{T_p}\to H_p$ is an isomorphism.

Choose $j_0\in U$.  Since $\ell_{j_0}(p)\ne0$, projection
$H_p\to\CC^{U\setminus\{j_0\}}$ is an isomorphism.  Work in the projective
chart $z_{j_0}\ne0$ and normalize $z_{j_0}=1$.  Choose a basis
$\lambda^{(1)},\ldots,\lambda^{(k-3)}$ of $W$ whose first $r-2$ vectors span
$W_{T_p}$.  For the remaining $k-r-1$ equations
\[
        G_a(z)=\sum_i\lambda_i^{(a)}z_i^n \qquad (a=r-1,\ldots,k-3),
\]
the Jacobian with respect to the $k-r-1$ unit variables
$z_j$, $j\in U\setminus\{j_0\}$, at $q$ is
\[
 \bigl(\lambda_j^{(a)}\bigr)_{a,j}
 \operatorname{diag}\bigl(nq_j^{\,n-1}\bigr)_{j\in U\setminus\{j_0\}}.
\]
The first factor is invertible by the preceding identification with $H_p$, and
the diagonal factor is invertible because every $q_j$, $j\in U$, is nonzero.
The formal implicit-function theorem therefore eliminates precisely these
$k-r-1$ unit variables as formal power series in the vanishing coordinates
$z_i$, $i\in T_p$.  After subtracting their values at $q$, the eliminated unit
coordinates have zero constant term, hence lie in the ideal generated by the
vanishing coordinates.  Consequently the completed maximal ideal is carried
\emph{exactly} to the vertex ideal $(x_i:i\in T_p)$, not merely to an ideal with
the same radical.

Crucially, the first $r-2$ equations are supported on $T_p$ and are not altered by
this elimination.  Thus the vanishing coordinates themselves survive as the
cone coordinates, and we obtain the coordinate-preserving isomorphism
\eqref{eq:completed-local-cone}, with
$F_a=\sum_{i\in T}\lambda_i^{(a)}x_i^n$ for $a=1,\ldots,r-2$.
For $r=2$ there are no remaining equations, so the completed local ring is
regular.

For $r\ge3$, we also record explicitly that the projective curve $C_{n,r}$ is
integral.  Write $u_1,\ldots,u_r$ for the restrictions to the corresponding
$\PP^1$ of the $r$ arrangement forms through $p$.  Their zeroes are pairwise
distinct.  In $\CC(\PP^1)^\times/\CC(\PP^1)^{\times n}$ the classes
\[
        u_2/u_1,\ldots,u_r/u_1
\]
are independent: valuation at the zero of $u_j$ detects the exponent of
$u_j/u_1$.  Hence the generic Kummer algebra is a field, and the associated
function-field extension has degree $n^{r-1}$.  On the other hand, the restriction
$C_{n,r}\to\PP^1$ of the coordinatewise power map has degree
\[
        \deg\OO_{C_{n,r}}(n)=n\deg C_{n,r}=n^{r-1}.
\]
Because the map to $\PP^1$ is finite, every irreducible component would dominate
$\PP^1$; the generic fiber is the spectrum of the field above, so there is a
unique minimal point and its generic local ring is a field.  Since $C_{n,r}$ is a
complete intersection, hence Cohen--Macaulay and $S_1$,
Lemma~\ref{lem:s1-generic-field} shows that $C_{n,r}$ is integral.  It is smooth because the line $\PP^1$ meets each
coordinate hyperplane in a distinct point, the cover is etale
away from those points, and locally at each branch point it is obtained by
adjoining an $n$-th root of a parameter.  Since $C_{n,r}$ is a complete
intersection of type $(n,\ldots,n)$ with $r-2$ factors in $\PP^{r-1}$, its degree
is $n^{r-2}$ and adjunction gives
\[
 \omega_{C_{n,r}}\simeq
 \OO_{C_{n,r}}((r-2)n-r),
\]
which yields \eqref{eq:genus}.

The cone ring in \eqref{eq:completed-local-cone} is a standard graded normal
domain: it is a complete intersection, hence $S_2$, and it is regular away from
the vertex because $C_{n,r}$ is smooth; the vertex has codimension two.  Moreover
the coordinate-preserving statement above shows
that the maximal ideal of $\OO_{X_n,q}$ becomes the vertex ideal after
completion.  Lemma~\ref{lem:completed-cone-blowup} therefore applies.  It shows
that the ordinary blow-up of $q$ is smooth and that its exceptional fiber is a
reduced copy of $C_{n,r}$, with
\[
        \OO(-E)|_{C_{n,r}}\simeq\OO_{C_{n,r}}(1).
\]
Hence its self-intersection is
$-\deg\OO_{C_{n,r}}(1)=-n^{r-2}$.
\end{proof}

\end{document}
